\sectionguard
\section*{The Propers: Themes and Movement}
\label{triptych:brief-synthesis:start}

\movementthesis{Four of this formulary's ten texts stop exactly where a human
contribution would begin. The Collect asks God to add what prayer does not
presume to ask; the Epistle's last words are grace not made void, and the
missal cuts before the Apostle's labour; the Communion antiphon drops the
imperative \latin{da ei} and keeps the filling of the barns; and in the Gospel
the man who is healed says nothing at all until the string of his tongue has
been loosed by somebody else. What holds the day together is not a theme but a
grammar: God acts first, and the texts decline to describe what we add.}

The book prints the ten texts in one order and the argument can be read in
four stages:

\begin{enumerate}[leftmargin=1.5em, itemsep=0.1em]
  \item A holy place that turns out to be a people, and a body that flowers again. \cue{Int.}, \cue{Grad.}, \cue{All.}
  \item What exceeds merit, and what exceeds request --- with the lesson cut one clause short. \cue{Coll.}, \cue{Ep.}
  \item \latin{Ephphetha}: one word, and a mouth that will not close. \cue{Gosp.}
  \item The healed man's psalm, the accepted gift, the overflowing press. \cue{Off.}, \cue{Sec.}, \cue{Comm.}, \cue{Postcomm.}
\end{enumerate}

English quoted in this section is the Douay--Rheims (Challoner) for scriptural
propers and the 1861 Cummiskey hand missal for orations, as throughout; the
witnesses are identified at each element in the commentary and in the
references.

\subsection*{1. A holy place that is a people, and a body that flowers again}

The Introit is built, not quoted. The missal's own reference --- \latin{6-7 et
36} --- announces that three clauses have been cut from three places in Psalm
67 and welded: God in his holy place; God who makes the like-minded dwell in a
house; he will give power and strength to his people. Read in that order the
antiphon performs a small argument. It names a place, and then says the place
is inhabited, and then says the strength belongs to the inhabitants. St.\
Augustine put the point sharply on the second clause: God is contained by no
place, so we should not look for him outside ourselves, but rather, dwelling
in the house of one mind, deserve that he should dwell in us --- and that
place, he adds, grace builds, not merits.

One word of the Introit is a small marvel of transmission, and it is argued at
the element below. The chant sings \latin{unánimes}, ``of one mind,'' where
every Bible of the Latin west reads \latin{unius moris}, ``of one manner.''
Augustine did not have the chant's word; he supplied it as his gloss. Peter
Lombard passed the gloss on; Bellarmine read it back into the verse.
Cassiodorus alone lemmatises it as text. The Introit sings what began as an
explanation.

The Gradual and the Alleluia carry that forward in two different registers.
The Gradual is Psalm 27:7 with v.~1 for its versicle --- thanksgiving first,
the earlier cry second, exactly the inversion the Introit performs on its own
psalm --- and its central words are \latin{reflóruit caro mea}, my flesh has
flowered again. Augustine glosses that in four words: \latin{id est, et
resurrexit caro mea}. Cassiodorus gives the prefix a history (the flesh
flowered in the virginal conception and re-flowered in the Resurrection);
Aquinas gives it a span from Adam to Christ; Bellarmine gives it a face
(``restored \dots\ to the bloom of youth, health, joy, and beauty''). Theodoret
declines the whole reading and hears a hunted man's recovered vigour. The day
sings the verse without settling the question. The Alleluia then does what a
festal chant does: it summons the singing --- take a pleasant psalm with the
harp --- and it does so immediately before the Gospel in which a man is given
back the power of speech.

\subsection*{2. What exceeds merit, and what exceeds request}

The Collect is two clauses of address and two of petition, and the two halves
are exactly matched. God is addressed as exceeding both \latin{mérita} and
\latin{vota} --- what suppliants deserve and what suppliants ask --- out of the
abundance of his kindness. Then: pardon what conscience fears, and add what
prayer does not presume. Below what we deserve, God does not exact; above what
we ask, God supplies. The Latin makes \latin{orátio} itself the thing that does
not presume, which is a stranger and better sentence than the 1861 hand
missal's ``we dare not presume to ask''; the divergence is declared at the
element. St.~Thomas states the principle without commenting on the prayer:
\latin{Deus multa nobis praestat ex sua liberalitate etiam non petita}.

The Epistle sets the same asymmetry in the first person. Paul delivers what he
received --- died, buried, raised, seen --- and then puts himself at the end of
the witness list as the untimely born, the least of the apostles, not worthy of
the name. And then the lesson's last words: \latin{grátia autem Dei sum id quod
sum, et grátia eius in me vácua non fuit}. That is where the missal stops, and
where it stops is the single most consequential fact about this reading. The
verse continues, in the Bible, ``but I have laboured more abundantly than all
they. Yet not I, but the grace of God with me'' --- and every classical
exposition of Pauline grace built on this text turns on those unread words.
Augustine's proof of free will is the labour; Gregory's
prevenient-and-subsequent scheme hangs on the single word \latin{mecum}; Aquinas derives
\latin{gratia cooperans} from the same clause. The cut is old --- the same
mid-verse ending stands in the Missals of 1570 and 1862 --- and its motive is
undocumented. Two things follow. A guide may not say that the Fathers expound
this lesson's teaching on grace: they expound one clause further. And the
truncation cannot be pressed into a denial of cooperation, since the very
authors whose key clause is missing are the ones who teach it.

\subsection*{3. \latin{Ephphetha}: one word, and a mouth that will not close}

The Gospel is a single healing told almost entirely in gestures. Jesus is in
Gentile country throughout --- Tyre, Sidon, the Decapolis --- and the man is
brought to him by others. He is taken out of the crowd; fingers are put in his
ears; there is spittle, and a touch on the tongue; there is a look upward and a
groan; and then one Aramaic word, which the evangelist translates:
\latin{Ephphetha, quod est adaperíre}. St.~Bede reads the fingers as the Holy
Spirit, proving it from Luke and Matthew on the finger of God, and the spittle
as divine wisdom instructing the mouths of those being catechised; he reads the
sigh not as need but as example. The seventh-century commentary passing under
Jerome's name reads the whole itinerary as salvation crossing from Judaea to
the Gentiles and names what the loosed mouth first says: \latin{Credo in Deum
Patrem omnipotentem} --- and the missal, at the end of this Gospel and no
other in the formulary, prints the rubric \latin{Credo}.

Then the closing paradox: he charged them to tell no one, and the more he
charged them the more they published it. St.~Augustine treats only this, and
takes it as a rebuke to lazy preachers. And the word itself walked out of the
Gospel into the sacraments: St.~Ambrose says plainly that the Church's rite of
opening comes from this passage, and the \work{Rituale Romanum} was still
printing \latin{Ephpheta, quod est, Adaperire} over the ears of infants in the
century that produced the missal.

\subsection*{4. The healed man's psalm, the accepted gift, the overflowing press}

Immediately after that Gospel the Offertory sings \latin{clamávi ad te, et
sanásti me} --- I cried to thee, and thou hast healed me. It is the day's one
proper spoken by somebody already healed, and it arrives at the moment the
gifts are set out. A caution belongs here rather than in an appendix, because
it changes the claim: the entire checked tradition, Latin and Greek alike,
reads that healing of the Resurrection, on the strength of the psalm's title,
the dedication of David's house. Jerome states it in one line; Augustine asks
who healed a man who was never sick, and answers that the flesh was healed when
it rose. To hear the verse as the deaf-mute's thanksgiving is to read it
against the tradition, not with it.

The Secret then asks one gift to do two things --- be acceptable to God, and be
a support to the frailty of those offering it --- and the Communion, uniquely
in this formulary, quotes not a psalm but the wisdom of Israel: honour the Lord
with thy substance and with the firstfruits, and thy barns shall be filled and
thy presses run over with wine. Bede reads the firstfruits at three levels, and
the middle one is the day's own Epistle in other words: a man honours God with
his substance who credits every good he has to grace and not to his own powers.
The Postcommunion closes by refusing to divide the beneficiary: help of mind
\textit{and} body, and, saved in both, \latin{gloriémur} --- may we glory --- in
the fullness of the heavenly remedy. Schuster draws the consequence that the
body's contact with the Body of Jesus is a pledge of the final resurrection,
and with that the day's last word rejoins the first: power and strength given
to his people.

\label{triptych:brief-synthesis:end}
\clearpage
